% Folder 136, batch 14 (archivists' pages 261-280).
% Pass: Opus 5.5 (claude-opus-5-5), 2026-10-03 — first pass, unchecked against the pages by a human.
%
% Pages 261, 263, 265, 267, 269, 271, 273, 275, 277, 279 are typescript versos
% unrelated to the notes and are skipped. Page numbers follow the batch order
% (page 1 of the batch = 261), which the scan footers confirm; from page 274
% on, the pencilled numbers on the leaves read one higher (275, 277, 279, 281).
\documentclass[11pt,a4paper]{article}
\input{../preamble/grothendieck}

\folder{136}
\batch{14}
\pages{261}{280}
\foldertitle{Complexe de De Rham à puissance divisée [conférence de 1976 à l’IHÉS] : notes manuscrites (s.d.).}
\dating{[à partir de 1975-1976]}
\watermark{Édition de démonstration}

\begin{document}

\page{262}
\note{le calcul qui ouvre la page continue un exemple commencé avant ce lot.}

\emph{Ex.} $M = S[-i]$, \ill{} d\ill{} …
\[
\begin{aligned}
H_i\bigl(R^{*}\varphi_N\,\tau_N(P)\bigr) &\simeq \mathbb{R}\operatorname{Hom}\bigl(\tau_N(S[-i]),\,P\bigr)\\
&\simeq \mathbb{R}\operatorname{Hom}\bigl((\tau_{N-i}S)[-i],\,P\bigr)\\
&\simeq \mathbb{R}\operatorname{Hom}^{i}(\tau_{N-i}S,\,P),
\end{aligned}
\]
donc
\[
\bigl(R^{n}\varphi_N(\tau_N P)\bigr)_i = \operatorname{Ext}^{n}\bigl((\tau_{N-i}S)[-i],\,P\bigr).
\]

\emph{NB} (c'est valable aussi en remplaçant $\varphi_N$ par $\varphi_{N',N}$ ($N'\leq N$), pourvu que $i \geq N'$ (si $i < N'$, on trouve $0$ bien sûr\,!)).

\struck{\emph{NB} Supposons $P \in \operatorname{Im\,ess}\varphi_{N',N}$ $(N'\leq N)$, i.e. $P \xrightarrow{\sim} \varphi_{N',N}\tau_N P = R^{0}\varphi_{N',N}(\tau_N P)$. Alors $R^{n}(\varphi_{N',N}(\tau_N P)) = 0$ si $n>0$, i.e. $\operatorname{Ext}^{n}(\tau_{N-i}(S)[-i],\,P) = 0$ si $i \geq N'$, $n>0$. En effet, cela revient à}
\note{passage encadré et barré de traits obliques ; suit une formule biffée, \struck{$R^{n}\varphi_N\tau(P)$}, avec un $M_N$ au-dessous.}

Soit $C$ une catégorie abélienne, \add{et $N'_0 \leq N$ dans $\mathbb{Z}$,} et
\[
C \xrightarrow{\;i\;} M_{N'_0}
\]
un foncteur \emph{exact}, tel que
\[
i(C) \subset \operatorname{Im\,ess}\varphi_{N'_0,N},
\]
i.e.
\[
i(M) \xrightarrow{\sim} \varphi_{N'_0,N}\,\tau_{N,N'_0}\,i(M).
\]
\note{l'indice que nous rendons par $N'_0$ porte un signe en exposant (prime ou astérisque) mal formé ; même lecture dans toute la page.}
Je dis que l'on a
\[
\bigl(R^{n}\varphi_{N'_0,N}\bigr)(\tau_{N,N'_0}\circ i) = 0 \quad\text{si } n>0.
\]
En effet, comme $i \simeq \varphi_{N'_0,N}\,\tau_{N,N'_0}\,i$, on conclut
\[
\begin{aligned}
R^{*}i &\simeq R^{*}\bigl(\varphi_{N'_0,N}\,\tau_{N,N'_0}\,i\bigr)\\
&\simeq \bigl(R^{*}\varphi_{N'_0,N}\bigr)(\tau_{N,N'_0}\circ i)
\quad\text{car } \tau_{N,N'_0}\circ i \text{ exacte},
\end{aligned}
\]
et $R^{*}i \simeq i$ car $i$ exact.

\page{264}
d'où la conclusion.

Le cas qui nous intéresse est
$N'_0 = 0 \leq N$, $C = \operatorname{Mod}(k)$, $i(M) = M\otimes_k S$, et on sait que
\[
M\otimes_k S \subset \operatorname{Im\,ess}\varphi_{0N} \qquad \forall M \in \operatorname{Mod}(k),
\]
i.e.
\[
M\otimes_k S \simeq \varphi_{0N}\,\tau_N(M\otimes_k S)
\quad\bigl(= R^{0}\varphi_{0N}(\tau_N M\otimes_k S)\bigr) ;
\]
alors on a
\[
R^{n}\varphi_{0N}\bigl(\tau_N(M\otimes_k S)\bigr) = 0 \quad\text{si } n>0,
\]
\note{au-dessous de $\varphi_{0N}\tau_N(M\otimes_k S)$, un signe $=$ vertical mène à une formule biffée, \struck{$\operatorname{Hom}((\tau_{N-i}S)[-i],$} \struck{$M\otimes$}.}
i.e. (encadré)
\[
\boxed{\operatorname{Ext}^{n}_{\mathcal{M}}\bigl((\tau_{N-i}S)[-i],\,M\otimes_k S\bigr) = 0 \quad\text{si } n>0,\ i\geq 0.}
\]

\[
\operatorname{Hom}_{\mathcal{M}}(M,\,\varphi_\omega P) \simeq \dots \simeq \varinjlim_N \operatorname{Hom}^{*}_{\mathcal{M}}(\tau_{N-i}M,\,P)
\]
\note{l'indice $N-i$ de $\tau$ est douteux à la première ligne. Sous le premier $\simeq$, un renvoi : « si $M$ de \uncertain{prés. finie} ».}

Si de plus $M$ projectif,
\[
\operatorname{Hom}\bigl(M,\,R^{*}\varphi_\omega(P)\bigr) \simeq \varinjlim_N \mathbb{R}\operatorname{Hom}_{\mathcal{M}}\bigl((\tau_{N-i}M)[-i],\,-\bigr),
\]
\note{un symbole est noirci avant la virgule dans le dernier membre.}
en particulier
\[
\operatorname{Hom}\bigl(M,\,R^{n}\varphi_\omega P\bigr) \simeq \varinjlim_N \operatorname{Ext}^{n}_{\mathcal{M}}\bigl((\tau_{N-i}M)[-i],\,P\bigr).
\]
P.\,ex. si $M = S[i]$, on trouve
\[
\begin{aligned}
\bigl(\text{\struck{$R^{*}$}}(\varphi_\omega P)\bigr)_i &= \varinjlim_N \operatorname{Hom}^{i}(\tau_{N-i}S,\,P),\\
\bigl(R^{*}\varphi_\omega(\tau_\omega P)\bigr)_i &\simeq \varinjlim_N R^{*}\operatorname{Hom}^{i}\bigl(\tau_{N-i}(S),\,P\bigr),\\
\bigl(R^{n}\varphi_\omega(\tau_\omega P)\bigr)_i &\simeq \varinjlim_N \operatorname{Ext}^{n}_{\mathcal{M}}\bigl(\tau_{N-i}(S)[-i],\,P\bigr).
\end{aligned}
\]
\note{la dernière ligne est précédée d'un début biffé illisible, \struck{\ill{}}.}

\page{266}
Si $C \xrightarrow{\;i\;} \mathcal{M}_{N_0}$ est un foncteur exact,
\[
\operatorname{Im\,ess} i \subset \operatorname{Im\,ess}\varphi_{N_0,\omega},
\]
alors on a $i \simeq \varphi_{N_0,\omega}\,\tau_{\omega,N_0}\,i$, d'où
\[
\bigl(R^{*}\varphi_{N_0,\omega}\bigr)(\tau_{\omega,N_0}\circ i) \simeq R^{*}i ;
\]
en particulier
\[
\bigl(R^{n}\varphi_{N_0,\omega}\bigr)\bigl(\tau_{\omega,N_0}\,i(M)\bigr) = 0 \quad\text{si } n>0,\ M\in\operatorname{ob}C.
\]
On est intéressé au même cas particulier qu'avant.

Soit $K^{\bullet}$ un complexe de $C$, $L^{\bullet}$ un complexe de $\mathcal{M}$\ill{}, et supposons donné un hom dans la catégorie dérivée $D(\mathcal{M}_{N_0})$
\[
i(K^{\bullet}) \xrightarrow{\;\alpha\;} L^{\bullet}
\]
tel que
\[
\tau_{\omega,N_0}\,i(K) \xrightarrow[\sim]{\;\tau_\omega\alpha\;} \tau_{\omega,N_0}L^{\bullet}
\]
soit un isom dans $D(\mathcal{M}_\omega)$, i.e. tel que $\forall n$,
\[
H^{n}(\alpha) : H^{n}\bigl(i(K)\bigr) \longrightarrow H^{n}(L^{\bullet}),
\qquad H^{n}\bigl(i(K)\bigr) \simeq i\bigl(H^{n}(K)\bigr),
\]
soit un \uncertain{quasi-isom}, i.e. noyau et conoyau sont négligeables. On en conclut

\page{268}
\uncertain{on obtient que} $\mathbb{R}\varphi_{0,\omega}$
\[
\bigl(\mathbb{R}^{*}\varphi_{N_0,\omega}\bigr)(\tau_{\omega,N_0}\,i)(K) \simeq \bigl(\mathbb{R}^{*}\varphi_{N_0,\omega}\bigr)(\tau_{\omega,N_0}L^{\bullet})
\]
i.e.
\[
i(K) \simeq \underbrace{\bigl(\mathbb{R}^{*}\varphi_{N_0,\omega}\bigr)(\tau_{\omega,N_0}L^{\bullet})}_{L'^{\bullet}}.
\]
\note{le membre de gauche de la première ligne est souligné d'un trait renvoyant à $i(K)$ de la seconde.}

\struck{Je dis que l'hom d'augmentation}
Mais c'est là un isom dans $D(\mathcal{M}_{N_0})$ seulement, — rien de plus précis que le domaine de départ — on aimerait un isom dans $D(C)$. Comme on suppose $i$ pl.\ fidèle et avec un adjoint à droite $k$, donc
\[
k\circ i \simeq \mathrm{id}_C,
\]
et on suppose $k$ \emph{exact}. Appliquant $k$ à l'isom obtenu, on trouve
\[
\begin{aligned}
K^{\bullet} &\simeq k\Bigl(\mathbb{R}^{*}\varphi_{N_0,\omega}\bigl(\tau_{\omega,N_0}L\bigr)\Bigr)\\
&\simeq \bigl(\mathbb{R}^{*}(k\,\varphi_{N_0,\omega})\bigr)(\tau_{\omega,N_0}L^{\bullet})
\end{aligned}
\]
i.e. $K^{\bullet}$ se \uncertain{reconstitue} dans $D(C)$ à l'aide de
\[
\tau_{\omega,N_0}\bigl(i(K^{\bullet})\bigr) \simeq \tau_{\omega,N_0}(L^{\bullet})
\]
dans $D(\mathcal{M}_\omega)$.
\struck{On va simplifier l'expression obtenue, en considérant}
\[
\text{\struck{$\varphi_{N_0,\omega}(\tau_{\omega,N_0}L^{\bullet}) \longrightarrow \mathbb{R}^{*}\varphi_{N_0,\omega}(\tau_{\omega,N_0}L)$}}
\]
\struck{\emph{Lemme.} Cet hom est un quasi-isomorphisme.}

\struck{\emph{Corollaire.}}
\[
\text{\struck{$K^{\bullet} \simeq k\,\varphi_{N_0,\omega}(\tau_{\omega,N_0}L^{\bullet})$}}
\]
\[
\text{\struck{$\simeq k\,\varphi_{N_0,\omega}(\tau_{\omega,N_0}\,i\,K)$}}
\]
\note{tout ce dernier passage, depuis « On va simplifier », est encadré et barré de traits obliques.}

\page{270}
\[
C \longrightarrow \mathcal{M}_\omega
\]
\ill{} $= \sigma$ ;\quad en retour : $k\,\varphi_{N_0,\omega} = \psi$.
\note{une flèche courbe revient de $\mathcal{M}_\omega$ vers $C$, marquée $k\varphi_{N_0\omega}=\psi$ ; le nom du foncteur aller, au-dessus, est surchargé et illisible (il s'agit vraisemblablement du composé de $i$ et de la troncation $\tau_{\omega,N_0}$).}

$\sigma$ pl.\ fidèle, adj.\ à droite $\psi$,
\[
\psi\sigma \xleftarrow{\;\sim\;} \mathrm{id}_C
\]
\struck{$\psi$ exact} ; $\sigma$ exact ; $k$ exact ; mais $\psi$ peut-être pas exact.

\struck{$\mathbb{R}$} $\sigma(K^{\bullet}) \longrightarrow L^{\bullet}_\omega$ quasi-isom dans $\mathcal{M}_\omega$.
\[
(R^{*}\psi)\,\sigma(K^{\bullet}) \xrightarrow{\;\sim\;} R^{*}\psi(L^{\bullet}_\omega) \xleftarrow[\;\sim\;]{?} \psi(L^{\bullet}_\omega)
\]
où $\sigma(K^{\bullet}) = R^{*}\sigma\,K$, et
\[
(R^{*}\psi)\,\sigma(K^{\bullet}) = R^{*}(\psi\sigma)(K) \simeq K \qquad (\psi\sigma \simeq \mathrm{id}).
\]
\[
H^{*}\bigl(\mathbb{R}\psi(L^{\bullet}_\omega)\bigr) \Longleftarrow E_2^{pq} = R^{q}\psi\bigl(H^{p}(L^{\bullet}_\omega)\bigr)
\]
si $R^{q}\psi\bigl(H^{p}(L^{\bullet}_\omega)\bigr) = 0$ pour $q>0$
\[
\text{\struck{$E_1^{p,q} \to E_1^{p+1,q}$,\quad $E_2^{p,q} \to E^{p+2,q-1}$}}
\]
On trouve
\[
H^{n}\bigl(\mathbb{R}\psi(L^{\bullet}_\omega)\bigr) \simeq \psi\bigl(H^{n}(L^{\bullet}_\omega)\bigr).
\]
\note{ce calcul par la suite spectrale est traversé de deux grands traits obliques ; le $R^{*}\psi$ placé au-dessus de $H^{*}$ est biffé.}

\[
\begin{aligned}
i(K^{\bullet}) &\longleftarrow L^{\bullet}\\
\tau_{\omega N_0}\,i(K^{\bullet}) &\rightleftarrows \tau_{\omega N_0}(L^{\bullet})\\
K^{\bullet} &\simeq \mathbb{R}^{*}(k\circ\varphi_\omega)\bigl(\tau_{\omega N_0}(L^{\bullet})\bigr)
\end{aligned}
\]
\note{ces trois lignes sont également traversées de traits obliques.}

\[
D\bigl(\operatorname{Mod}(k)\bigr) \longrightarrow D(\mathcal{M}_\omega) \quad\text{pl.\ fidèle}
\]

\page{272}
\struck{$\psi\sigma = \mathrm{id}$}

\begin{tikzcd}
D^{+}(C) \arrow[r, bend left=20, "R\sigma"] & D^{+}(\mathcal{M}_\omega) \arrow[l, bend left=20, "R\psi"]
\end{tikzcd}

\[
\psi\sigma \xleftarrow{\;\sim\;} \mathrm{id}_C, \qquad \sigma\psi \longrightarrow \mathrm{id}_{\mathcal{M}_\omega}
\]
\note{ces deux relations sont écrites en haut à droite de la page.}

\begin{tikzcd}[column sep=small, row sep=normal, nodes={font=\small}]
 & & {\operatorname{Hom}(R\sigma K,\, R\sigma R\psi(L))} \arrow[dll] \\
{\operatorname{Hom}(R\sigma(K),\, L)} \arrow[rr, "{?\ \sim}"] \arrow[dr] & & {\operatorname{Hom}(K,\, R\psi(L))} \arrow[u] \arrow[dl, no head, "\wr" description] \\
 & {\operatorname{Hom}(R\psi R\sigma K,\, R\psi L)} &
\end{tikzcd}

avec $R\sigma R\psi(L) = \mathbb{R}(\sigma\psi)(L)$ et $R\psi R\sigma K = R(\psi\sigma)(K) \simeq K$.

\begin{tikzcd}
{R\sigma R\psi R\sigma K} \arrow[d] & {R\sigma R\psi L} \arrow[d] \\
R\sigma K & L
\end{tikzcd}

\[
D^{+}C \longrightarrow D^{+}\mathcal{M}_\omega
\]
pl.\ fid., compatible avec multiplication.

\page{274}
\[
C \xrightarrow{\;i\;} \mathcal{M}, \qquad M \longmapsto M\otimes_k S
\]
\[
i \xrightarrow{\;\sim\;} \varphi_{0,N}\,\tau_{N,0}\, i
\]
\note{un symbole est noirci entre $\tau_{N,0}$ et $i$.}
\[
M \hookrightarrow M^{[N,+\infty[} \rightrightarrows M^{[N,+\infty]\times\mathbb{N}}
\]
\[
x \longmapsto (c_{i,k-i}\,x)_{k\geq N}, \qquad
(x_k)_{k\geq N} \rightrightarrows \bigl(c_{\ell,k}\,x_k - c_{\ell,k-i}\,x_{k+\ell}\bigr)
\quad (0\leq i\leq N)
\]
\note{les deux flèches parallèles sont celles d'un noyau ; l'indice $0\leq i\leq N$ est écrit au-dessus de la formule de droite.}
\[
M \longmapsto \underline{\operatorname{Hom}}^{*}_{k}(S,\,M)
\]
\[
\varphi^{0,N}\,\tau^{N,0}\circ i \longrightarrow i \qquad
\text{syst.\ des } \underline{\operatorname{Hom}}_k(S_i,\,M) \simeq S\otimes_k M
\]
\note{au-dessous : « $\mathrm{Ab}$ » et « $\operatorname{Ind}(\mathrm{Ab})$ ».}
\[
M\otimes_k S_i \longrightarrow \underline{\operatorname{Hom}}^{i}\bigl(\tau_{N-i}(S),\,M\otimes_k S\bigr)
\]

$C$ cat.\ ab.\ $k$-linéaire (i.e.\ $k \to \operatorname{End}(\mathrm{id}_C)$),

$X$ objet de $C$,

$M$ un $k$-module \ill{} (\ill{} dans

$\underline{\operatorname{Hom}}(M,X) \in C$ défini par
\[
\operatorname{Hom}\bigl(Y,\,\underline{\operatorname{Hom}}(M,X)\bigr) = \operatorname{Hom}_k\bigl(M,\,\operatorname{Hom}_C(Y,X)\bigr),
\]
$M\otimes_k X$ défini par
\[
\operatorname{Hom}(M\otimes_k X,\,Y) \simeq \operatorname{Hom}_k\bigl(M,\,\operatorname{Hom}_C(X,Y)\bigr).
\]
Dans le \ill{} \uncertain{ça existe} ;
\[
M \longmapsto \underline{\operatorname{Hom}}(M,X) \quad\text{et}\quad \text{\struck{$M\mapsto$}}\ M \longmapsto M\otimes_k X
\]
\uncertain{est exact à droite} —

\page{276}
\[
P \to Q \to R, \qquad P \to Q \to Z' \to 0, \qquad Z' \to R
\]
\[
\text{\struck{$P \to Z \to 0$\quad $0 \to Z \to Q \to Z' \to 0$}}
\]
\uncertain{considérons} $T$ ;
\[
0 \to Z' \to R, \qquad T(P) \to T(Q) \to T(Z') \to 0
\]

\begin{tikzcd}[column sep=small]
T(Z) \arrow[r] & T(Q) \\
T(P) \arrow[u] &
\end{tikzcd}

\[
T(Q)/I \simeq T(Z') \overset{?}{\hookrightarrow} T(R), \qquad I = \operatorname{Im}\bigl(T(P)\bigr)
\]
\[
0 \to Z \to Q \to Z' \hookrightarrow R, \qquad T(Z) \to T(Q) \to T(Z') \to 0
\]
\[
I \hookrightarrow R
\]
\uncertain{vérifier} \ill{} \uncertain{injective en} \ill{}

\uncertain{transformer en surjection} \uncertain{par} $\operatorname{Hom}$
\[
Z' \hookrightarrow R
\]
\[
P \to Q \to I \hookrightarrow R, \qquad I \to 0
\]
\[
T(P) \leftarrow T(Q) \leftarrow T^{\bullet}(I) \leftarrow 0, \qquad
T(I) = \operatorname{Ker}\bigl(T(Q) \to T(P)\bigr)
\]
\[
\operatorname{Hom}(P,M), \qquad \operatorname{Hom}_{C,S}\bigl(\mathbb{N},\,\operatorname{Hom}(P,M)\bigr)
\]
\[
\text{\struck{$\operatorname{Hom}(\mathbb{N}\otimes P,\,M)$}}, \qquad \operatorname{Hom}_{\otimes}
\]
\note{une grande ligne courbe sépare, sur la page, les calculs de gauche de cette colonne de droite.}
\[
\begin{matrix}
T(R) & & & \\
\text{\struck{$\mathbb{Z}^{[N,+\infty[\times\mathbb{N}}$}} & \mathbb{Z}^{[N,+\infty]} & & \mathbb{Z}\\
\| & \| & & \|\\
P & \longrightarrow\ Q & \longrightarrow & R \longrightarrow 0\\
 & & & \cup\\
 & & & R'
\end{matrix}
\]
\note{une flèche oblique va de $Q$ à $R'$.}
\[
(\lambda_{k,\ell})_{k\geq N,\ \ell\geq 0} \longmapsto \sum_{k,\ell}\lambda_{k,\ell}\bigl(c_{\ell,k}\,e_k - c_{\ell,k-i}\,e_{k+\ell}\bigr)
\]
\[
(\lambda_k)_{k\geq 0} \longmapsto \sum \ill{}
\]
\note{la dernière expression est entourée et entièrement hachurée ; on y devine $c_{\ell,k-i}\lambda_k$.}

\page{278}
\[
S = k[U], \qquad U^{n} = U_n \text{ de degré } -n
\]
ou $S^{k} = k\{\{T\}\}$ \ill{} \uncertain{à priori} en $S$-modules (\uncertain{gradués}) avec $S = k\{T\}$.

$M \in \operatorname{Mod}(k)$,

$M \longmapsto M\otimes_k S$ un foncteur
\[
C = \operatorname{Mod}(k) \xrightarrow{\;i_*\;} \operatorname{Mod\,gr.}\,k\{T\} = \mathcal{M}.
\]
Si $N\in\mathbb{Z}$, $\sigma_N(P) = P/\tau_{N+1}(P)$ « \uncertain{cotronqué} » \uncertain{en} \uncertain{dim} $N$ ($\leftarrow$ « \uncertain{tronqué à gauche} »),
\[
\sigma_N(P)_i = \begin{cases} P_i & \text{si } i\leq N,\\ 0 & \text{si } i>N.\end{cases}
\]
Considérons le composé $\sigma_N\circ i_*$ ($N\leq 0$, on fera $N\to-\infty$).

\emph{Th.} Le foncteur $\tau_N\,i_*$ est pl.\ fid.
\[
\operatorname{Hom}_{\text{\struck{gr}}\,\mathcal{M}}\bigl(M\otimes_k S,\,\tau^{N}P\bigr) = \operatorname{Hom}_k\bigl(M,\,\operatorname{Hom}_{\mathcal{M}}(S,\,\ill{})\bigr)
\]
\[
\operatorname{Hom}_{\mathcal{M}_N}\bigl(\tau^{N}M\otimes_k S,\,\tau^{N}N\otimes_k S\bigr) = \operatorname{Hom}_{\mathcal{M}}\bigl(M\otimes_k S,\,i^{N}\tau^{N}(N\otimes_k S)\bigr)
\]
\[
= \operatorname{Hom}_k\bigl(M,\,\operatorname{Hom}_{\mathcal{M}}(S,\,i^{N}\tau^{N}-)\bigr)
\]
\[
\operatorname{Hom}_{\mathcal{M}}\bigl(S,\,\tau^{N}N\otimes_k S\bigr) \simeq N\ ??
\]
\note{au-dessous, en regard : $\operatorname{Hom}_{V(S),\,\mathrm{gr}}\bigl(V(N\otimes S),\,V(S)\bigr)$.}
\[
\operatorname{Hom}_S(S,\,M\otimes_k S), \qquad \operatorname{Hom}_{\mathrm{coalg}}(S,\,M\otimes_k S)
\]
\[
= \operatorname{Hom}_{V(S)\text{-}\ill{}}\bigl(\operatorname{Hom}_k(M,\,V(S)),\,V(S)\bigr)
\]
\[
\text{\struck{$D$}}\ V_S(M) = V(M\otimes_k S) = \bigl(k' \mapsto \operatorname{Hom}_k(M,\,S^{k'})\bigr)
\]
$V(S)$ \uncertain{monoïde}, \ \ $\operatorname{Hom}$ \ill{} $V_S(M)$.

À gauche de la page :
\[
U_{n+k} \xrightarrow[\;T^{(k)}\;]{\;c_{n,k}\;} U_n, \qquad \operatorname{P}_\infty
\]
\[
\eta_n \in P_{-n}, \qquad \xi_n\,T^{(k)}\,\eta_{n+k} = c_{n,k}\,\eta_n
\]
\[
\text{\struck{\ill{}}}, \qquad \eta_{n+k} = \xi_n \ill{}
\]
\[
\xi_{n+k}\,T^{(k)}\,U_{n+k} = \xi_n\,c_{n,k}\,U_n, \qquad T^{(k)}U_{n+k} = c_{n,k}\,U_n,
\]
\[
c_{n,k}\,(\eta_{n+k} - \xi_n) = 0.
\]
\note{le placement des expressions de cette page, en deux colonnes séparées par un trait, est simplifié ; la lecture des indices de $\eta$ et $\xi$ est incertaine.}

\page{280}
\note{la moitié supérieure de la page porte les calculs ci-dessous, à l'endroit ; la moitié inférieure, écrite tête-bêche, est donnée ensuite.}
\[
M\otimes S\otimes M'\otimes S, \qquad M\otimes M'\otimes S\otimes S
\]
\[
\begin{aligned}
&\mathrm{id}_M\otimes u_S\otimes \mathrm{id}_{M'}\otimes \mathrm{id}_S, &\quad &\mathrm{id}_{M'}\otimes u_S\otimes \mathrm{id}_S,\\
&\mathrm{id}_M\otimes \mathrm{id}_S\otimes \mathrm{id}_{M'}\otimes u_S, & &\mathrm{id}_{M'}\otimes \mathrm{id}_S\otimes u_S
\end{aligned}
\]
\note{ces quatre lignes, d'une encre plus pâle, sont écrites dans la marge du haut.}
\[
S\boxtimes S = S, \qquad S\boxtimes C \simeq C, \qquad SL
\]
\[
\boxed{(M\otimes_k S)\boxtimes_S C \simeq M\otimes_k C}\ ?
\]
\[
\{M\otimes_k S\}\otimes_k C \longleftarrow M\otimes_k C, \qquad
\begin{aligned}
&\mathrm{id}_M\otimes u_S\otimes \mathrm{id}_C\\
&\mathrm{id}_M\otimes \mathrm{id}_S\otimes u_C
\end{aligned}
\]
\note{en regard, à droite, « $\mathrm{MON}$ » et un petit cercle pointé, $\odot$.}
\[
C \longrightarrow C\otimes_k S \rightrightarrows C\otimes_k S\otimes_k S
\]
\[
\text{\struck{$V(S)\otimes V(C) \to V(C')$}}, \qquad
\operatorname{Hom}_k(M,\,S^{k'})\times C^{k'} \longrightarrow C'^{k'}
\]
\[
V(M\otimes_k S)(k') = \operatorname{Hom}_k(M,\,S^{k'}),
\]
\[
\operatorname{Hom}_{k'}(M\otimes_k S,\,k') = \operatorname{Hom}_k(M,\,S^{k'}) \quad\text{est un } S^{k'}\text{-module\,!}
\]

\note{moitié inférieure, tête-bêche ; elle est traversée d'un long trait oblique et, en partie, entourée.}
\[
\alpha_*(\mathbb{1}),\ \gamma_0^{(i)}, \qquad
\Delta\gamma_0^{i} = \sum_{j+k=i} \gamma_0^{j}\otimes\gamma_0^{k}
\]
\[
\beta_*(t) = \sum t^{(i)}\gamma_0^{i} = \exp(t\gamma_0)
\]
\[
A^{n}\,\widehat{\otimes}\,A^{n} \longleftarrow A^{n}
\]
\[
\bigl(A^{*}(\Sigma,\Sigma_+,t)\bigr)^{k'} = A_*\bigl(\Sigma^{k'},\Sigma_+^{k'},t^{k'}\bigr),
\qquad A_*(S,S^{+},t)
\]
\[
A^{*}(\Sigma,\Sigma_+,t)^{\vee} = A_*^{\wedge}\bigl(\check\Sigma,\check\Sigma_+,\check t\bigr)
\]
$\Sigma/k$ : $\Sigma \supset \Sigma'$ \uncertain{coalgèbre} avec \struck{\ill{}} \uncertain{sous-coalgèbre} \struck{\ill{}} \uncertain{et puissances divisées} \ill{}, \ill{} \uncertain{on se donne} $\Sigma \xrightarrow{\;t\;} k$ \ill{} $\Sigma'$.
\[
0 \longrightarrow \text{\struck{$\Sigma'$}}^{\,k+} \longrightarrow \Sigma^{k} \longrightarrow \Sigma'^{k}, \qquad t_k, \qquad k\{\{T\}\}\{\{\gamma_0,\,-\gamma_0\}\}
\]
\note{sous la suite exacte, « l'idéal » ; le passage entouré, sur les coalgèbres à puissances divisées et leur sous-coalgèbre, est très raturé et nos lectures y restent douteuses. Le calcul se poursuit vraisemblablement au-delà du lot.}

\end{document}
